\documentclass{article}
\usepackage{amsmath,amssymb}
\usepackage[margin=1in]{geometry}

\newcommand{\CC}{\mathcal{C}}
\newcommand{\NN}{\mathcal{N}}
\newcommand{\union}{\cup}

\begin{document}

\section*{Clade-level NNI expansion of a CCD}

\paragraph{Motivation.}
The CCD0 construction \emph{expands the split set}: for every clade partition
$P\to\{L,R\}$ that could exist because its parent $P$ and both children $L,R$ are
already clades in $\CC$ (the set of clades with nonzero marginal support), it
adds that partition even if no sampled tree realised it. Here we describe the
analogous operation one level up (expanding the \emph{clade set} $\CC$ itself)
by reading nearest-neighbour-interchange (NNI) rearrangements directly off
the stored splits, without ever building a tree.

\paragraph{The move at the clade level.}
An internal edge of a tree is a stored split $P\to\{L,R\}$: the parent clade $P$
branches into disjoint children $L,R$ with $L\union R=P$. An NNI at that edge
needs the grandchildren. Descend into one child, say $L$, with its split
$L\to\{A,B\}$, and let the sibling be $R$. The three subtrees meeting at the edge
are $A$, $B$, $R$. The two NNI moves recombine them so that the child of $P$ is no
longer $L=A\union B$ but
\begin{equation}
  A\union R \qquad\text{or}\qquad B\union R .
  \label{eq:recomb}
\end{equation}
Because $A\subseteq L$ and $L\cap R=\varnothing$, each union in
\eqref{eq:recomb} is a disjoint union of size $\ge 2$ and a \emph{proper}
subclade of $P$ (it omits the nonempty $B$ resp.\ $A$). Each novel clade carries
an NNI-implied \emph{seed split}
\[
  (A\union R)\to\{A,R\},
\]
whose two constituents $A,R$ are themselves clades already in $\CC$. The roles of
$L$ and $R$ are symmetric, so a split $R\to\{C,D\}$ likewise yields $C\union L$
and $D\union L$. Every quantity ($A,B,R$ and the parent $P$) is read
straight from the CCD's split table.

\paragraph{Algorithm.}
For each clade $P$ with $|P|\ge 3$, each split $\{L,R\}\in\Pi(P)$, each choice of
descend-child $X\in\{L,R\}$ with sibling $Y$, and each grandchild split
$\{A,B\}\in\Pi(X)$, emit the candidates $A\union Y$ and $B\union Y$; keep those
not already in $\CC$. This is a pass over pairs of adjacent splits, costing
\[
  O\!\Big(\textstyle\sum_{P} |\Pi(P)|\cdot \max_{X}|\Pi(X)|\Big)
  \;=\; O\big(|\text{splits}|\times \bar{d}\big),
\]
with $\bar d$ the average number of splits per clade, proportional to the size
of the CCD, not to the number of trees it represents.

\paragraph{Two pairing modes.}
The CCD graph mixes splits from many posterior trees, so the algorithm above can
combine a split $P\to\{L,R\}$ with a grandchild split $L\to\{A,B\}$ that never
occurred together in any one tree. We therefore distinguish:
\begin{itemize}
  \item \textbf{Graph-wide} ($\NN_{\mathrm{gw}}$): pair \emph{any} stored split on
    $P$ with \emph{any} stored split on its child $X$. Cheapest and broadest; the
    pass described above.
  \item \textbf{Co-occurring} ($\NN_{\mathrm{co}}$): only recombine grandchild
    splits that actually appeared together under $P$ in some single base tree,
    the true NNI neighbourhood of the posterior sample. This is obtained by
    walking the stored base trees and reading $A,B,R$ off each internal edge.
\end{itemize}
By construction $\NN_{\mathrm{co}}\subseteq\NN_{\mathrm{gw}}$: the co-occurring
set is the graph-wide set restricted to split pairs realised within a single
tree. The graph-wide extras are exactly the clades that \emph{no single sampled
tree is one NNI move away from}; whether that is desirable broader coverage or an
overreach beyond the sample's NNI neighbourhood is a modelling choice one should
state explicitly.

\paragraph{Conservativeness.}
NNI is the most local rearrangement, so the expansion recovers only clades
``near'' the support and misses anything requiring an SPR-scale move. For the
single tree $(((A,B),C),D)$ the expansion yields exactly $\{ABD, CD, AC, BC\}$;
it does \emph{not} produce $AD$, $BD$, $ACD$ or $BCD$, which need a non-local
move. If the concern is held-out trees that differ from the sample by a single
local rearrangement, this coverage is exactly right; for broader topological
novelty one would extend to SPR neighbourhoods, also expressible at the clade
level (a regrafted subtree clade unioned with whatever lies below the target
edge), at the cost of many more candidates.

\paragraph{Empirical behaviour.}
On a real posterior (RSV2, $n=129$ taxa, $2251$ trees after burn-in), $\CC$ has
$1331$ clades. The graph-wide expansion finds $|\NN_{\mathrm{gw}}|=21{,}198$
novel clades (a $15.9\times$ blow-up); the co-occurring expansion finds
$|\NN_{\mathrm{co}}|=4{,}253$ ($3.2\times$). The containment
$\NN_{\mathrm{co}}\subseteq\NN_{\mathrm{gw}}$ holds with no co-occurring-only
clades, and the cross-tree recombinations account for the $\approx 5\times$ gap.
The whole pass runs in a few seconds.

\section*{Parameterisation in CCD1 / regCCD}

The clade expansion is most natural in the \emph{split}-parameterised models.
A regularised CCD (regCCD) takes the CCD1 graph, adds the unobserved splits whose
parent and children are already clades (the CCD0 split expansion), and prices
every split of a clade $C$ by additive smoothing,
\begin{equation}
  \theta_\alpha(S)=\frac{f(S)+\alpha}{f(C)+\alpha\,|S(C)|},
  \label{eq:smoothing}
\end{equation}
a $\mathrm{Dirichlet}(\alpha,\dots,\alpha)$ prior over the splits $S(C)$. This
gives nonzero probability to trees with \emph{unsampled splits of sampled
clades}, but a tree containing an \emph{unsampled clade} still has probability
zero: the residual limitation of both CCD0 and CCD1.

NNI clade expansion closes this gap, and, crucially, needs \emph{no new
parameter}. Under the CCD1 form $P(T)=\prod_v \mathrm{CCP}\!\left(\text{split}\mid
\text{clade}\right)$ a clade carries no marginal parameter; its probability is
induced by the product along the path to it. So adding a novel clade $K=A\cup R$
only requires its \emph{splits} to be in the graph:
\begin{itemize}
  \item the recombined parent split $P\to\{K,B\}$ enlarges $S(P)$ and is priced
    by \eqref{eq:smoothing} with $f=0$, i.e.\ mass $\alpha/(f(P)+\alpha|S(P)|)$;
  \item the seed split $K\to\{A,R\}$ (and any further splits of $K$ into existing
    clades found by the split expansion) is priced the same way; a clade left
    with a single split is deterministic ($\mathrm{CCP}=1$).
\end{itemize}
A held-out tree that is one NNI move from the support is then scored
$P(T)=\bigl[\prod\text{observed CCPs}\bigr]\cdot\theta_\alpha(P\to\{K,B\})\cdot
\theta(K\to\{A,R\})>0$. The novel clades enter the \emph{same} additive-$\alpha$
Dirichlet--multinomial as the observed and split-expanded splits (single
pseudocount), so the only model change is a larger clade set.

\paragraph{Coverage vs.\ sharpness, and a two-level prior.}
The benefit is out-of-sample coverage of trees with unsampled clades; the cost is
that enlarging $|S(P)|$ dilutes the smoothed mass of observed splits. The latter
is controlled by giving NNI-derived splits (those touching a novel clade) a
\emph{separate} pseudocount $\beta$, while observed and split-expanded splits keep
$\alpha$:
\begin{equation}
  \theta(S)=\frac{f(S)+\alpha\,[S\text{ regular}]+\beta\,[S\text{ NNI}]}
                 {f(C)+\alpha\,|R(C)|+\beta\,|N(C)|},
  \label{eq:twolevel}
\end{equation}
an asymmetric Dirichlet placing a weaker prior on the speculative
recombinations. Crucially $\beta$ does not affect \emph{coverage} (any
$\beta>0$ keeps the NNI split positive); it only rebalances mass at observed
parents.

We evaluate by \emph{batched leave-one-tree-out}: $B$-fold cross-validation that
rebuilds each model on the training folds and scores the held-out fold, so every
tree is scored once out-of-training. A full rebuild per fold is needed because
the NNI clade set depends on which trees are present, so the count-subtraction
shortcut used for plain regCCD does not apply; batching makes this affordable
($B$ rebuilds rather than $N$). To track true leave-one-tree-out under an
autocorrelated chain we assign trees to folds in a \emph{strided} fashion (tree
$i$ to fold $i \bmod B$), which keeps each held-out tree's neighbours in training
as LOTO does; this converges to LOTO as $B\to N$.

On an RSV2 posterior ($n=129$, $2251$ trees, $\alpha=0.4$, strided) we report
$B=20$ folds; the estimate is converged there (coverage is identical at $B=40$,
and within one tree of $B=10$), so it approximates leave-one-tree-out. The
fraction of held-out trees with nonzero probability rises from $90.6\%$ (regCCD)
to $99.5\%$ (co-occurring) and $100.0\%$ (graph-wide). With a single pseudocount
($\beta=\alpha$) the mean log probability on the commonly covered trees drops
from $-53.5$ to $-55.5$ (co-occurring), the dilution cost. Shrinking $\beta$
recovers it at unchanged coverage: $\beta=0.10$ gives $-54.1$ and $\beta=0.02$
gives $-53.6$, essentially matching the baseline while still covering $99.5\%$ of
held-out trees. So the
conservative co-occurring mode with a small $\beta$ buys the coverage gain at
negligible sharpness cost. A single held-out tree ($1/2251$) is never covered,
even by graph-wide: it needs more than one NNI move (an SPR-scale rearrangement).

\paragraph{Choosing the pseudocounts.}
Since $\beta$ does not change which trees are covered, the held-out log
probability over the (fixed) covered set is a clean objective for selecting
$(\alpha,\beta)$, optimised by the same batched cross-validation. Because the
pseudocounts do not change the graph, each fold is built once and the grid is
scored in closed form, so the search costs only the $B$ fold builds. On RSV2 this
selects, for the co-occurring mode, $\alpha\approx 0.2$ and $\beta\approx 0.01$:
the optimal $\beta$ is an order of magnitude below $\alpha$, confirming the weak
prior on NNI splits, and the optimal $\alpha$ is below the value used without
clade expansion, so the pseudocounts should be re-tuned rather than inherited.

\end{document}
